\documentclass[11pt,a4paper]{article}
\usepackage{fontspec,amsmath,amssymb,amsthm,microtype}
\setmainfont{lmroman10-regular.otf}[BoldFont=lmroman10-bold.otf,ItalicFont=lmroman10-italic.otf,BoldItalicFont=lmroman10-bolditalic.otf]
\setmonofont{lmmono10-regular.otf}[Scale=MatchLowercase]
\usepackage[margin=25mm]{geometry}
\usepackage[unicode,hidelinks]{hyperref}
\setlength{\emergencystretch}{3em}
\setlength{\parskip}{0.4em}
\setlength{\parindent}{0pt}
\allowdisplaybreaks[2]
\raggedbottom
\title{\texorpdfstring{The genuine quadratic Taylor remainder}{The genuine quadratic Taylor remainder}}
\author{LangevinLean mathematical exposition}
\date{}
\begin{document}\maketitle
\noindent Source module: \texttt{Kinetic\allowbreak{}Langevin/\allowbreak{}Hessian\allowbreak{}Remainder.\allowbreak{}lean}. The accompanying HTML pairs this mathematical exposition with the exact Lean statements and records the scope of the formalization.


The second-order integrator linearizes the force at its initial position. Its error requires a quadratic estimate for this actual linearization, obtained from the explicitly integrated Hessian along a segment. The derivative and Hessian Lipschitz assumptions are ordinary analytic hypotheses on the force.


\section{\texorpdfstring{Integrate Hessian variation along the segment}{Integrate Hessian variation along the segment}}


The averaged Hessian is the Bochner integral of H(y+t(x\ensuremath{-}y)), with t between zero and one. Subtracting H(y) inside the integral gives an integrand with norm at most M\ensuremath{{}_{2}}t\ensuremath{\Vert}x\ensuremath{-}y\ensuremath{\Vert}. The norm of an integral is at most the integral of this nonnegative envelope, whose value is M\ensuremath{{}_{2}}\ensuremath{\Vert}x\ensuremath{-}y\ensuremath{\Vert}/\allowbreak{}2. Lean checks continuity of the integrand, moves subtraction through the integral, applies the Lipschitz inequality on the segment, and evaluates the elementary integral of t.



{\raggedright\noindent Lean declarations: \texttt{averaged\allowbreak{}Hessian\_\allowbreak{}sub\_\allowbreak{}base\_\allowbreak{}bound}.\par}

\section{\texorpdfstring{Apply the averaged operator to the displacement}{Apply the averaged operator to the displacement}}


The fundamental theorem of calculus identifies g(x)\ensuremath{-}g(y) with the averaged Hessian applied to x\ensuremath{-}y. Thus the true Taylor remainder equals (\ensuremath{\bar{H}}\ensuremath{-}H(y))(x\ensuremath{-}y). The operator norm inequality bounds this by \ensuremath{\Vert}\ensuremath{\bar{H}}\ensuremath{-}H(y)\ensuremath{\Vert}\ensuremath{\Vert}x\ensuremath{-}y\ensuremath{\Vert}, yielding M\ensuremath{{}_{2}}\ensuremath{\Vert}x\ensuremath{-}y\ensuremath{\Vert}\ensuremath{{}^{2}}/\allowbreak{}2. This is an actual remainder estimate; the Hessian is never replaced by an assumed scalar or a fictitious common mean-value point.



{\raggedright\noindent Lean declarations: \texttt{force\_\allowbreak{}taylor\_\allowbreak{}remainder\_\allowbreak{}bound}.\par}

\end{document}
